\documentclass[11pt]{article}
\usepackage[margin=1in]{geometry}
\usepackage{parskip}
\usepackage{tocloft}
\renewcommand{\cftsecpresnum}{}
\renewcommand{\cftsecaftersnum}{.}
\usepackage{titlesec}
\titleformat{\section}{\normalfont\Large\bfseries}{\thesection.}{1em}{}
\titleformat{\subsection}{\normalfont\bfseries}{\thesubsection.}{1em}{}
\usepackage{amssymb,amsmath,amsfonts,verbatim,amsthm, amscd}
\usepackage[dvipsnames]{xcolor}
\usepackage[breakable,skins]{tcolorbox}
\usepackage{graphicx,tikz-cd,adjustbox}
\usepackage{float}
\usepackage{scalerel}
\usepackage{stackengine}

\newtcolorbox{mybox}[2][]{
    arc=0mm, enhanced, frame hidden, breakable
}
\usepackage{mathtools}

\newcommand{\nocontentsline}[3]{}
\newcommand{\tocless}[2]{\bgroup\let\addcontentsline=\nocontentsline#1{#2}\egroup}

\usepackage{subcaption}
\usepackage{array} % core package
\usepackage[sc]{mathpazo} % Palatino font with small caps
\usepackage[scr=euler]{mathalfa}
\usepackage[unicode=true]{hyperref} % permits navigating on the pdf
\usepackage{enumerate}
\hypersetup{
    colorlinks=true,
    linkcolor=NavyBlue,
    citecolor=OliveGreen,
    urlcolor=blue,
    pdfencoding=auto
}

\linespread{1.1}
\setlength\parindent{0.2in}
\usepackage{cleveref} % smart referencing
\usepackage{caption}

% --- Math Commands from User ---
\newcommand{\calL}{\mathcal{L}}
\newcommand{\calH}{\mathcal{H}}
\newcommand{\bz}{\mathbb{Z}}
\newcommand{\lt}{\normalfont\text{lt}}
\newcommand{\calA}{\mathcal{A}}
\newcommand{\calS}{\mathcal{S}}
\newcommand{\ind}{\text{\normalfont Ind}}
\newcommand{\stab}{\text{\normalfont Stab}}
\newcommand{\chow}{\normalfont\text{\underline{H}}}
\newcommand{\augchow}{\normalfont\text{H}}
\newcommand{\klpoly}{\normalfont\text{P}}
\newcommand{\invklpoly}{\normalfont\hat{\text{Q}}}
\newcommand{\redchar}{\overline{\chi}}
\newcommand{\rank}{\normalfont\text{rk}}
\newcommand{\corank}{\normalfont\text{crk}}
\newcommand{\flats}{\normalfont\text{Flats}}
\newcommand{\flags}{\normalfont\text{Flags}}
\newcommand{\matM}{\normalfont\text{M}}
\newcommand{\hilbMmodF}{\normalfont\text{\underline{H}}_\text{M/F}}
\newcommand{\M}{\normalfont\text{M}}
\newcommand{\N}{\normalfont\text{N}}
\newcommand{\rk}{\normalfont\text{rk}}
\newcommand{\chibar}{\overline{\chi}}
\newcommand{\E}{\normalfont\text{E}}
\newcommand{\ZZ}{\mathbb{Z}}
\newcommand{\Ebasis}[1]{\mathrm{#1}}
\newcommand{\Cbasis}[1]{\mathrm{c}^{#1}}
\newcommand{\CbasisSub}[2]{\mathrm{c}^{#1}_{#2}}
\newcommand{\klbasis}[1]{\mathrm{z}^{#1}}
\newcommand{\klbasisSub}[2]{\mathrm{z}^{#1}_{#2}}
\newcommand{\Ibasis}[2]{\mathfrak{b}^{#1}_{#2}}
\newcommand{\pal}{\normalfont\text{Pal}}
\newcommand{\klchowcoeff}{\mathfrak{a}}
\newcommand{\intM}{\normalfont\text{Int}}
\newcommand{\mathscrI}{\mathscr{I}}
\newcommand{\calP}{\mathcal{P}}
\newcommand{\IntP}{\normalfont\text{Int}}
\newcommand{\rev}{\normalfont\text{rev}}
\newcommand{\qbinom}[2]{\begin{bmatrix} #1 \\ #2 \end{bmatrix}_q}

% --- Theorem Environments ---
\newtheorem{theorem}{Theorem}
\newtheorem{corollary}[theorem]{Corollary}
\newtheorem{lemma}[theorem]{Lemma}
\newtheorem{proposition}[theorem]{Proposition}
\newtheorem{conjecture}[theorem]{Conjecture}
\newtheorem{example}[]{Example}
\newtheorem{definition}[theorem]{Definition}

\theoremstyle{remark}
\newtheorem*{remark}{Remark} % Unnumbered remark

\title{Notes: The $\pal(0)$-module of flats}
\author{Nutan Nepal}
\date{\today}

\begin{document}
\maketitle

\section{Algebraic Preliminaries}

\subsection{Posets and Incidence Algebras}
Let \(\calP\) be a \emph{bounded, weakly ranked, partially ordered set}.
We denote by \(\IntP(\calP)\) the set of all closed intervals of \(\calP\) and
by \(\rho: \IntP(\calP) \to \ZZ\)
the \emph{weak rank function}.

The \emph{incidence algebra} of \(\calP\) over \(\ZZ[x]\), denoted \(\mathscr{I}\), is the free \(\ZZ[x]\)-module on \(\IntP(\calP)\). An element \(a \in \mathscr{I}\) is a function that assigns a polynomial \(a_{st}(x)\) to each interval \([s,t]\). The algebra structure is given by the convolution product:
\[ (ab)_{su} := \sum_{s\leq t\leq u} a_{st}(x) \cdot b_{tu}(x), \qquad
    \text{ for every \(s\leq u\) in \(\calP\).}
\]
This algebra is associative and has a multiplicative identity \(\delta\).
An element \(a \in \mathscr{I}\) is invertible if and only if \(a_{ss}(x)\) is a unit
for every \(s \in \calP\). Two fundamental elements in \(\mathscr{I}\) are the \emph{zeta function},
\(\zeta_{st}(x) = 1\), and its inverse, the \emph{Möbius function}, \(\mu = \zeta^{-1}\).

\subsection{Kernels, KLS functions, and the Chow function}
The weak rank function \(\rho\) induces a subalgebra \(\mathscr{I}_\rho \subseteq \mathscr{I}\) consisting of elements \(a\) such that \(\deg(a_{st}) \le \rho_{st}\) for every interval \([s,t]\). This subalgebra is equipped with an involution \(a \mapsto a^{\rev}\) defined by:
\[ (a^{\rev})_{st}(x) := x^{\rho_{st}} a_{st}(x^{-1}) \]
An element \(\kappa \in \mathscr{I}_\rho\) is a \emph{\((\calP,\rho)\)-kernel} if
\(\kappa_{ss}(x) = 1\) and it satisfies the property \(\kappa^{-1} = \kappa^{\rev}\).
The primary example is the \emph{characteristic function} \(\chi := \mu \cdot \zeta^{\rev}\),
which is a kernel for any weakly ranked poset.

In \cite{ferroni-matherne-vecchi}, the authors show that a kernel uniquely determines two key functions:
\begin{theorem}\label{thm:kernel-properties}
    Let \(\kappa \in \mathscrI_\rho\) be a \((\calP,\rho)\)-kernel. There exist unique elements
    \(f, g \in \mathscrI_\rho\) (the right and left KLS functions, respectively) satisfying:
    \begin{enumerate}[\normalfont(i)]
        \item \(f_{ss}(x) = 1\) and \(g_{ss}(x) = 1\) for all \(s \in \calP\).
        \item \(\deg f_{st}(x) < \frac{1}{2} \rho_{st}\) and \(\deg g_{st}(x) < \frac{1}{2}
              \rho_{st}\) for all \(s < t\).
        \item \(f^{\rev} = \kappa \cdot f\) and \(g^{\rev} = g \cdot \kappa\).
    \end{enumerate}
\end{theorem}

From a kernel \(\kappa\) and its KLS functions, one can define several other
important polynomial invariants.
Central to this paper is the \emph{Chow function},
which is formulated in \cite{ferroni-matherne-vecchi} as follows:

\begin{definition}\label{def:chow-reduced} [Reduced Kernel and Chow Function]
    Let \(\kappa\) be a \((\calP,\rho)\)-kernel.
    \begin{enumerate}
        \item The \emph{reduced kernel} \(\overline{\kappa} \in \mathscr{I}_\rho\) is given by:
              \[ \overline{\kappa}_{st}(x) = \begin{cases} \frac{\kappa_{st}(x)}{x - 1} & \text{if } s < t, \\ -1 & \text{if } s = t.\end{cases}\]
        \item The \emph{\(\kappa\)-Chow function} \(\chow \in \mathscr{I}_\rho\) is defined as:
              \[ \chow = - (\overline{\kappa})^{-1}\]
    \end{enumerate}
\end{definition}

\subsection{The Module of Flats and Perverse Subgroups}
Let \(R := \ZZ[x, x^{-1}]\) be the ring of Laurent polynomials with integer coefficients.
For any integer \(n\), we denote by \(\pal(n)\) the additive group of polynomials \(f(x) \in R\)
satisfying \(f(x) = x^n f(x^{-1})\). While \(\pal(0)\) is a subring of \(R\), we note that
\(\pal(n)\) for \(n \neq 0\) is not, as \(\pal(n) \cdot \pal(m) \subseteq \pal(n+m)\).

Let \(\M\) be a matroid with lattice of flats \(\mathcal{L}(\M)\).
\begin{definition}
    The module \(\calH(\M)\) is the free \(R\)-module with a standard basis \(\{\Ebasis{F}\}_{F \in \mathcal{L}(\M)}\). An element \(\alpha \in \calH(\M)\) is a formal sum
    \[ \alpha = \sum_{F \in \mathcal{L}(\M)} \alpha_F \cdot\Ebasis{F}, \quad \text{where } \alpha_F \in R. \]
\end{definition}

We also define the \(R\)-linear maps \(S_F: \calH(\M) \to R\) for any flat \(F \in \mathcal{L}(\M)\) by
\[ S_F(\alpha) := \sum_{G \ge F} x^{\rank(F)-\rank(G)} \alpha_G. \]
Following Braden-Vysogorets, there is an important abelian subgroup of \(\calH(\M)\).
\begin{definition}
    The subgroup of \textbf{perverse elements} \(\calH_p(\M) \subset \calH(\M)\) consists of all \(\alpha\) such that for every flat \(F \in \mathcal{L}(\M)\):
    \begin{enumerate}
        \item The coefficient \(\alpha_F\) is a polynomial, i.e., \(\alpha_F \in \ZZ[x]\).
        \item The Laurent polynomial \(S_F(\alpha)\) is in \(\pal(0)\).
    \end{enumerate}
\end{definition}
\begin{definition}
    We relax the first condition to define a larger subgroup \(\calH'_p(\M) \subseteq \calH(\M)\)
    consisting of all elements \(\alpha\) that satisfy \(S_F(\alpha) \in \pal(0)\)
    for all \(F \in \mathcal{L}(\M)\).
\end{definition}

\subsection{The Chow and Kazhdan-Lusztig Bases}
The module \(\calH(\M)\) has natural bases constructed from the Kazhdan-Lusztig and Chow polynomials.

\begin{definition}
    The \textbf{Kazhdan-Lusztig basis} \(\{\klbasis{F}\}\) is defined by
    \[ \klbasis{F} = \sum_{G \le F} \klbasisSub{F}{G}\cdot\Ebasis{G} := \sum_{G \le F} x^{\rank(F)-\rank(G)} \klpoly_{\M_G^F}(x^{-2}) \cdot \Ebasis{G}, \]
    where \(\klpoly_{\M}\) is the Kazhdan-Lusztig polynomial of \(\M\).
\end{definition}

\begin{remark}
    Braden-Vysogorets prove that the elements \(\{\klbasis{F}\}\) form a \(\ZZ\)-basis for \(\calH_p(\M)\). In
    particular, if \(\alpha = \sum_F \alpha_F \cdot \Ebasis{F} \in \calH_p(\M)\), then
    \[ \alpha = \sum_F \alpha_F(0) \cdot \klbasis{F}. \]
\end{remark}

\begin{definition}
    The \textbf{Chow basis} \(\{\Cbasis{F}\}\) is defined by
    \[
        \Cbasis{F} = \sum_{G\le F} \CbasisSub{F}{G}\cdot\Ebasis{G} :=
        \sum_{G \le F} x^{\rank(F)-\rank(G)} \chow_{\M_G^F}(x^{-2}) \cdot \Ebasis{G},
    \]
    where \(\chow_{\M}\) is the Chow polynomial of \(\M\), which is palindromic of degree \(\rank(\M)-1\).
\end{definition}

\begin{remark}
    The coefficients of the Chow basis have a special structure:
    \[\CbasisSub{F}{F} = 1 \in \pal(0),\quad\text{and}\quad\CbasisSub{F}{G} \in \pal(2)\ \text{for all}\ G < F.\]
\end{remark}
\begin{definition}
    We define the module \(\calS(\M)\) to be the free \(\pal(0)\)-module with basis \(\{\Cbasis{F}\}\).
\end{definition}

\begin{definition}
    For an element \(e\) in the ground set of \(\M\), the \textbf{deletion operator} is an \(R\)-module homomorphism \(\Delta : \calH(\M) \to \calH(\M \setminus e)\) defined by
    \[ \Delta(\Ebasis{F}) = x^{-\delta(F)}\cdot \Ebasis{F \setminus e}, \quad \text{where } \delta(F) := \rank_{\M}(F) - \rank_{\M\setminus e}(F \setminus e).\]
\end{definition}
\begin{remark}
    The deletion operators commute: \(\Delta_e \circ \Delta_f = \Delta_f \circ \Delta_e:
    \calH(\M) \to \calH(\M \setminus \{e,f\})\).
\end{remark}

\begin{definition}
    For any element \(e\) in the ground set of \(\M\), we define the set
    \[ S_e(\M) := \{ F \in \mathcal{L}(\M) \mid e \notin F \text{ and } F\cup e \in \mathcal{L}(\M) \}. \]
    As demonstrated in Braden-Vysogorets, \(S_e(\M)\) plays a crucial role in the deletion formula.
\end{definition}

\section{Main Propositions}

\begin{proposition}\label{prop:cinHp}
    For all \(F \in \mathcal{L}(\M)\), \(\Cbasis{F} \in \calH'_p(\M)\), i.e.,
    it satisfies \(S_G(\Cbasis{F}) \in \pal(0)\) for all \(G \in \mathcal{L}(\M)\).
\end{proposition}
\begin{proof}
    We see that, for any \(G \in \mathcal{L}(\M)\),
    \begin{align*}
        S_G(\Cbasis{F}) = \sum_{H: G \le H \le F} x^{\rank(G)-\rank(H)}\cdot\CbasisSub{F}{H}
         & = \sum_{H: G \le H \le F} x^{\rank(G)-\rank(H)} \left( x^{\rank(F)-\rank(H)} \chow_{\M_H^F}(x^{-2}) \right) \\
         & = x^{\rank(F)-\rank(G)} \sum_{H: G \le H \le F} x^{-2(\rank(H) -\rank(G))} \chow_{\M_H^F}(x^{-2}).
    \end{align*}
    This sum is precisely the definition of the \textit{augmented} Chow polynomial, \(\augchow_{\M_G^F}(x^{-2})\). So,
    \(S_G(c^F) = x^{\rank(F)-\rank(G)} \augchow_{\M_G^F}(x^{-2})\). Let \(d = \rank(F)-\rank(G)\).
    The term \(x^d\) is in \(\pal(2d)\).
    The term \(\augchow_{\M_G^F}(x^{-2})\) is in \(\pal(-2d)\).
    The product is in \(\pal(2d - 2d) = \pal(0)\) as required.
\end{proof}
\begin{remark}
    The proof of Proposition \ref{prop:cinHp} demonstrates that the Chow basis elements have the property that their image under the evaluation maps $S_G$ lies in $\pal(0)$ for all flats $G$.
\end{remark}
\begin{proposition}\label{prop:Hp_is_pal0_mod}
    The subgroup \(\calH'_p(\M)\) is a \(\pal(0)\)-module.
\end{proposition}
\begin{proof}
    Let \(\alpha = \sum_F p_F \cdot \Ebasis{F} \in \calH'_p(\M)\) and \(q \in \pal(0)\).

    We see that
    \[ S_G(q \cdot \alpha) = \sum_{H \ge G} x^{\rank(G)-\rank(H)} (q \cdot p_H)
        = q \cdot S_G(\alpha) \in \pal(0). \]
    Therefore, \(q \cdot \alpha \in \calH'_p(\M)\).
\end{proof}

\begin{proposition}\label{prop:LPDecomposition}
    Any Laurent polynomial \(f\) can be uniquely written as a sum \(f = f_0 + f_2\) where
    \(f_0 \in \pal(0)\) and \(f_2 \in \pal(2)\). The projections are given by
    \[ f_0(x) = \frac{x^2 f(x^{-1}) - f(x)}{x^2 - 1} \quad \text{and} \quad f_2(x) = \frac{f(x) - f(x^{-1})}{1 - x^{-2}}. \]
\end{proposition}

\begin{proposition}\label{prop:SisHp}
    \(\calS(\M) = \calH'_p(\M)\) as \(\pal(0)\)-modules.
\end{proposition}
\begin{proof}
    Proposition \ref{prop:cinHp} and \ref{prop:Hp_is_pal0_mod}
    show that \(\calS(\M) \subseteq \calH'_p(\M)\).
    It is now enough to show that \(\calH'_p(\M) \subseteq \calS(\M)\).

    Let \(\alpha = \sum_F \alpha_F \cdot \Ebasis{F}\) be an arbitrary element of \(\calH'_p(\M)\). Since \(\{\Cbasis{F}\}\) is a basis for the entire module \(\calH(\M)\), \(\alpha\) has a unique expansion in this basis:
    \[ \alpha = \sum_{F \in \calL(\M)} q_F \cdot \Cbasis{F}\quad \text{for } q_F \in \ZZ[x, x^{-1}]. \]
    To show that \(\alpha \in \calS(\M)\), we must prove that each of these coefficients \(q_F\) is actually in
    \(\pal(0)\). We prove this by induction on the corank of \(F\).

    \textbf{Base Case:} The relationship between the coefficients is given by the change-of-basis formula:
    \[ \alpha_F = \sum_{G \ge F} q_G \cdot \CbasisSub{G}{F} = q_F + \sum_{G > F} q_G \cdot \CbasisSub{G}{F}. \]
    For \(F=E\), this simplifies to \(q_E = \alpha_E\). Since \(\alpha \in \calH'_p(\M)\),
    we have \(S_E(\alpha) = \alpha_E \in \pal(0)\). This implies that \(q_E \in \pal(0)\), establishing the base case.

    \textbf{Inductive Step:} Assume that for all flats \(G\) with \(\rank(G) > \rank(F)\),
    the coefficient \(q_G\) is in \(\pal(0)\).

    Since \(\alpha \in \calH'_p(\M)\), it satisfies:
    \[ S_F(\alpha) = \sum_{H \ge F} x^{\rank(F)-\rank(H)} \alpha_H \in \pal(0). \]
    We substitute the change-of-basis formula \(\alpha_H = \sum_{L \ge H} q_L \cdot \CbasisSub{L}{H}\) into this expression:
    \[ S_F(\alpha) = \sum_{H \ge F} x^{\rank(F)-\rank(H)} \left( \sum_{L \ge H} q_L \cdot \CbasisSub{L}{H} \right)
        = \sum_{L \ge F} q_L \left( \sum_{H: F \le H \le L} x^{\rank(F)-\rank(H)} \CbasisSub{L}{H} \right)
        = \sum_{L \ge F} q_L \cdot S_F(\Cbasis{L}). \]
    We isolate the \(L=F\) term and rearrange the sum:
    \[ q_F = S_F(\alpha) - \sum_{L > F} q_L \cdot S_F(\Cbasis{L}). \]
    By the inductive hypothesis, each \(q_L\) in the sum is in \(\pal(0)\).
    Since \(\pal(0)\) is a ring, the difference on the right is in \(\pal(0)\).
    Thus, \(q_F\) must also be in \(\pal(0)\).
\end{proof}

\begin{proposition}[Change of Basis Formula]\label{prop:changeofbasis}
    If \(\alpha = \sum_F \alpha_F \cdot\Ebasis{F} \in \calS(\M)\),
    then its expansion in the \(\Cbasis{F}\) basis is
    \(\alpha = \sum_F \overline{\alpha_F} \cdot\Cbasis{F}\), where \(\overline{\alpha_F}\) is the unique \(\pal(0)\) component of \(\alpha_F\) from the decomposition in Proposition~\ref{prop:LPDecomposition}.
\end{proposition}
\begin{proof}
    Let \(\alpha \in \calS(\M)\). By the definition of \(\calS(\M)\), \(\alpha\) can be uniquely expressed as a \(\pal(0)\)-linear combination of the \(\Cbasis{F}\) basis vectors:
    \[ \alpha = \sum_{F \in \calL(\M)} q_F \cdot \Cbasis{F}, \quad \text{where each } q_F \in \pal(0). \]
    We can also express \(\alpha\) in the standard basis as \(\alpha = \sum_F \alpha_F \cdot\Ebasis{F}\).
    As in the proof of proposition \ref{prop:SisHp}, the change-of-basis formula relating the two sets of coefficients is:
    \begin{equation}\label{eq:changeofbasis}
        q_F = \alpha_F - \sum_{G > F} q_G \cdot \CbasisSub{G}{F}.
    \end{equation}
    The goal is to show that \(q_F = \overline{\alpha_F}\). We
    prove this by induction on the corank of \(F\) as in the previous proposition.

    \textbf{Base Case:} For the base case, we see that \(q_E = \alpha_E = S_E(\alpha) \in \pal(0)\).\
    This implies that \(q_E = \overline{\alpha_E}\).

    \textbf{Inductive Step:} Assume that for all flats \(G\) with \(\rank(G) > \rank(F)\),
    the coefficient \(q_G = \overline{\alpha_G}\).

    Then the entire sum \(\sum_{G > F} q_G \cdot \CbasisSub{G}{F}\) is in \(\pal(2)\), since for any \(G > F\),
    the coefficient \(q_G\) is in \(\pal(0)\) and the basis coefficient \(\CbasisSub{G}{F}\) is in
    \(\pal(2)\), as noted in a previous remark.
    This shows that \(\alpha_F\) is the sum of \(q_F\) (which, by definition, is in \(\pal(0)\))
    and another term that is in \(\pal(2)\).

    By Proposition~\ref{prop:LPDecomposition}, this decomposition is unique. Therefore,
    the coefficient \(q_F\) must be the unique \(\pal(0)\) component of the \(\alpha_F\).
\end{proof}


\begin{proposition}\label{prop:DeltaFInS}
    For any flat \(F \in \mathcal{L}(\M)\), the element \(\Delta(\Cbasis{F})\) is in \(\calS(\M \setminus e)\).
\end{proposition}
\begin{proof}
    Let \(F \in \mathcal{L}(\M)\) be an arbitrary flat. We need to show that
    \(S_H^{\M\setminus e}(\Delta(\Cbasis{F}))\) is in \(\pal(0)\) for all \(H \in \calL(\M \setminus e)\).

    Let \(\beta = \Delta(\Cbasis{F}) = \sum_{G \in \calL(\M \setminus e)} \beta_G \cdot \Ebasis{G}\).

    By definition, \(\Cbasis{F} = \sum_{L \le F} \CbasisSub{F}{L} \cdot \Ebasis{L}\). Applying the \(\Delta\) operator gives:
    \[ \Delta(\Cbasis{F}) = \sum_{L \le F} \CbasisSub{F}{L} x^{-\delta(L)}\cdot \Ebasis{L \setminus e}. \]
    The coefficient \(\beta_G\) is the sum over all flats \(L \le F\) such that \(L \setminus e = G\).
    Thus, we have:
    \begin{align*}
        S_H^{\M\setminus e}(\Delta(\Cbasis{F})) = \sum_{G \ge H} x^{\rank(H)-\rank(G)} \beta_G
         & = \sum_{G \ge H} x^{\rank(H)-\rank(G)} \left( \sum_{L: L \le F, L\setminus e = G} \CbasisSub{F}{L} x^{-\delta(L)} \right) \\
         & = \sum_{L: L \le F, L\setminus e \ge H} x^{\rank(H)-\rank(L \setminus e)} \CbasisSub{F}{L} x^{-\delta(L)}
    \end{align*}
    The exponent on \(x\) simplifies  to \( \rank(H) - \rank(L) \) using the definition of discrepancy.
    \[ S_H^{\M\setminus e}(\Delta(\Cbasis{F})) = \sum_{L: L \le F, L\setminus e \ge H} x^{\rank(H)-\rank(L)}
        \CbasisSub{F}{L}. \]
    Now we apply Lemma 2.15 from Braden-Vysogorets.
    The condition \(L\setminus e \ge H\) is equivalent to \(L \ge \overline{H}\), where \(\overline{H}\)
    is the closure of \(H\) in \(\M\). Furthermore, \(\rank_{\M\setminus e}(H) = \rank_{\M}(\overline{H})\).

    Substituting these two facts, the summation range becomes \(L \in [\overline{H}, F]\) and the expression is:
    \[ S_H^{\M\setminus e}(\Delta(\Cbasis{F})) = \sum_{L: \overline{H} \le L \le F} x^{\rank(\overline{H})-\rank(L)} \CbasisSub{F}{L}
        = S_{\overline{H}}(\Cbasis{F})\in \pal(0). \]
    Since the palindromic condition holds for any flat \(H\), we have proven that
    \(\Delta(\Cbasis{F}) \in \calS(\M \setminus e)\).
\end{proof}

\begin{remark}
    The map \(\Delta : \calS(\M) \to \calS(\M \setminus e)\) is a \(\pal(0)\)-module homomorphism that
    restricts from the \(\ZZ[x, x^{-1}]\)-module homomorphism \(\Delta : \calH(\M) \to \calH(\M \setminus e)\).
\end{remark}

\begin{proposition}\label{prop:imageofDeltaCbasis}
    For any flat \(F \in \mathcal{L}(\M)\), the following identities hold in \(\calS(\M \setminus e)\):
    \begin{enumerate}
        \item If \(e\) is not a coloop in \(\M^F\), then
              \begin{equation}\label{eq:imageofDeltaCbasisnoncoloop}
                  \Delta(\Cbasis{F}) = \Cbasis{F \setminus e} + x^{-1}\sum_{G \in S_e(\M^F)}\CbasisSub{F}{G\cup e} \cdot \Cbasis{G}
              \end{equation}
        \item If \(e\) is a coloop in \(\M^F\), then
              \begin{equation}\label{eq:imageofDeltaCbasiscoloop}
                  \Delta(\Cbasis{F}) = (x+x^{-1})\Cbasis{F \setminus e} + x^{-1}\sum_{G \in S_e(\M^F), G \neq F\setminus e}\CbasisSub{F}{G\cup e} \cdot \Cbasis{G}
              \end{equation}
    \end{enumerate}
\end{proposition}
\begin{proof}
    By Proposition~\ref{prop:changeofbasis}, the coefficients
    \(q_G\) in the expansion \(\Delta(\Cbasis{F}) = \sum q_G \cdot \Cbasis{G}\) are the unique \(\pal(0)\)
    components of the standard-basis coefficients,
    \[\beta_G = \sum_{L: L \le F, L\setminus e = G} \CbasisSub{F}{L} x^{-\delta(L)}.\]
    Thus we analyze
    \(q_G\) by considering different cases for the flat \(G\).

    \textbf{Coefficients for \(G \neq F \setminus e\):}
    For any flat \(G\) not equal to \(F \setminus e\), its preimages \(L \le F\) are a subset of
    \(\{\overline{G}, G \cup e\}\) where \(\overline{G}\) is the closure of \(G\) in \(\M\).
    \begin{itemize}
        \item If \(G \in S_e(\M^F)\), \(\overline{G} = G\) and the standard-basis coefficient is \(\beta_G = \CbasisSub{F}{G} + x^{-1}\CbasisSub{F}{G\cup e}\). This is a sum of a \(\pal(2)\) term and a \(\pal(0)\) term. Its unique \(\pal(0)\) part is \(q_G = x^{-1}\CbasisSub{F}{G\cup e}\).
        \item If \(G \notin S_e(\M^F)\), the only possible preimage is
              \(L=\overline{G}\). The standard-basis coefficient is
              \(\beta_G = \CbasisSub{F}{\overline{G}}\), which is in \(\pal(2)\). Its \(\pal(0)\) part is \(q_G = 0\).
    \end{itemize}
    These coefficients account for the summation in both formulas of the proposition.

    \textbf{Coefficient for \(G = F \setminus e\):}
    Here we analyze \(\beta_{F\setminus e}\) by considering the the following two cases:
    \begin{itemize}
        \item If \(e\) is \textbf{not a coloop} in \(\M^F\), the only flat \(L \le F\) that deletes to
              \(F \setminus e\) is \(L=F\). The standard-basis coefficient is \(\beta_{F\setminus e} =
              \CbasisSub{F}{F}x^{-0} = 1 \in \pal(0)\).
              So \(q_{F\setminus e} = 1\), proving the first part of the proposition.

        \item If \(e\) is \textbf{a coloop} in \(\M^F\), the preimages are \(L=F \setminus e\) and
              \(L=F\). The standard-basis coefficient is
              \(\beta_{F\setminus e} = \CbasisSub{F}{F\setminus e} + x^{-1}\CbasisSub{F}{F} = x+x^{-1}\in \pal(0)\).
              So \(q_{F\setminus e} = x+x^{-1}\). This proves the second part of the proposition.
    \end{itemize}
\end{proof}

\section{Recovery of the Semi-Small Decomposition Formulas}

We now show how the formulas in Proposition~\ref{prop:imageofDeltaCbasis} recover the known
semi-small decomposition formulas for the Chow polynomial \(\chow\) from [FMSV24]. This
is achieved by equating the coefficients of the empty flat \(\Ebasis{\emptyset}\) on both sides of the identities.

Let \(f = \rank(F)\). The coefficient of \(\Ebasis{\emptyset}\) in \(\Delta(\Cbasis{F})\) is:
\[ \beta_\emptyset = \CbasisSub{F}{\emptyset}x^{-\delta(\emptyset)} + \CbasisSub{F}{\{e\}}x^{-\delta(\{e\})} = x^f \chow_{\M^F}(x^{-2}) + x^{f-2}\chow_{\M^F_e}(x^{-2}) \]

\textbf{Case 1: \(e\) is not a coloop in \(\M^F\).}
We take the coefficient of \(\Ebasis{\emptyset}\) in the formula \eqref{eq:imageofDeltaCbasisnoncoloop}.
\[ x^{f}\chow_{\M^F}(x^{-2}) + x^{f-2}\chow_{\M^F_e}(x^{-2}) = x^{f}\chow_{\M^F\setminus e}(x^{-2}) +
    x^{-1}\sum_{G \in S_e(\M^F)}x^{f-g-1}\chow_{\M^F_{G\cup e}}(x^{-2})\cdot x^g\chow_{\M^G}(x^{-2}). \]
Dividing by \(x^f\) and setting \(y=x^{-2}\), we get:
\[ \chow_{\M^F}(y) + y\chow_{\M^F_e}(y) = \chow_{\M^F\setminus e}(y) + y\sum_{G \in S_e(\M^F)}\chow_{\M^F_{G\cup e}}(y)\cdot \chow_{\M^G}(y)\]
The sum on the right includes the term for \(G=\emptyset\), which is \(\chow_{\M^F_e}(y) \cdot \chow_{U_{0,0}}(y) = \chow_{\M^F_e}(y)\). Separating this term and canceling it with the \(y\chow_{\M^F_e}(y)\) from the LHS yields the semi-small decomposition formula:
\begin{tcolorbox}
    When \(e\) is not a coloop in \(\M^F\),
    \begin{equation} \label{eq:semi-small-decomposition-noncoloop}
        \chow_{\M^F}(y) = \chow_{\M^F\setminus e}(y) + \sum_{G \in S_e(\M^F)\setminus\{\emptyset\}} y \chow_{\M^F_{G\cup e}}(y) \cdot \chow_{\M^G}(y).
    \end{equation}
\end{tcolorbox}

\textbf{Case 2: \(e\) is a coloop in \(\M^F\).}
We take the coefficient of \(\Ebasis{\emptyset}\) in the formula \eqref{eq:imageofDeltaCbasiscoloop}.
Equating the coefficients gives:
\[ x^f \chow_{\M^F}(x^{-2}) + x^{f-2}\chow_{\M^F_e}(x^{-2}) = (x+x^{-1})x^{f-1}\chow_{\M^F\setminus e}(x^{-2}) + x^{f-2}\sum_{\dots} \dots \]
Dividing by \(x^f\) and setting \(y=x^{-2}\) leads to the known formula for adding a coloop:
\begin{tcolorbox}
    When \(e\) is a coloop in \(\M^F\), (some errors here?? looks good to me)
    \begin{equation} \label{eq:semi-small-decomposition-coloop}
        \chow_{\M^F}(y) = (1+y)\chow_{\M^F\setminus e}(y) +
        \sum_{G \in \mathcal{L}(\M^F\setminus e)\setminus\{\emptyset\}}y \chow_{\M^F_{G\cup e}}(y) \cdot
        \chow_{\M^G}(y).
    \end{equation}
\end{tcolorbox}
This completes the recovery of the known decomposition formulas of \(\chow\).

\subsection{Recovery of the known Augmented Chow Polynomial Decomposition Formulas}
We consider another \(\ZZ[x, x^{-1}]\)-module homomorphism
\(\Phi_{\M} : \calH(\M) \to \ZZ[x, x^{-1}]\) defined by
\[\Phi_{\M}(\alpha) = \sum_{F \in \mathcal{L}(\M)} x^{-\rank(F)} \alpha_F,\]
which satisfies \(\Phi_{\M\setminus e}\circ \Delta = \Phi_{\M}\).
Applying this map to \(\Cbasis{F}\) gives
\[\Phi_{\M}(\Cbasis{F}) = \sum_{G \le F} x^{-\rank(G)} \Cbasis{F}_G =
    x^{\rank(F)}\sum_{G \le F} x^{-2\rank(G)} \chow_{\M_G^F}(x^{-2}) =
    x^{\rank(F)} \augchow_{\M^F}(x^{-2}) \in \pal(0).\]
\begin{remark}
    We see that \(\Phi_{\M} = S^{\M}_{\emptyset}\). Thus,
    \(S^{\M\setminus e}_{\emptyset} \circ \Delta = S^{\M}_{\emptyset}\).
\end{remark}

Applying the \(\Phi_{\M\setminus e}\) map to \eqref{eq:imageofDeltaCbasisnoncoloop} yields the following identity:
\[ x^f \augchow_{\M^F}(x^{-2}) = x^f \augchow_{\M^F\setminus e}(x^{-2}) +
    x^{-1}\sum_{G \in S_e(\M^F)} x^{f-g-1}\chow_{\M^F_{G\cup e}}(x^{-2}) \cdot x^g\augchow_{\M^G}(x^{-2}). \]
Simplifying the exponents and setting \(y=x^{-2}\), we get:
\begin{tcolorbox}
    When \(e\) is not a coloop in \(\M^F\),
    \begin{equation} \label{eq:augchow-decomposition-noncoloop}
        \augchow_{\M^F}(y) = \augchow_{\M^F\setminus e}(y) +
        \sum_{G \in S_e(\M^F)} y \chow_{\M^F_{G\cup e}}(y) \cdot \augchow_{\M^G}(y).
    \end{equation}
\end{tcolorbox}
Proceeding similarly for the coloop case, we get the following identity:
\[ x^f \augchow_{\M^F}(x^{-2}) = (x+x^{-1})x^{f-1}\augchow_{\M^F\setminus e}(x^{-2}) +
    x^{-1}\sum_{G \in S_e(\M^F)} x^{f-g-1}\chow_{\M^F_{G\cup e}}(x^{-2}) \cdot x^g\augchow_{\M^G}(x^{-2}). \]
Simplifying the exponents and setting \(y=x^{-2}\), we get:
\begin{tcolorbox}
    When \(e\) is a coloop in \(\M^F\),
    \begin{equation} \label{eq:augchow-decomposition-coloop}
        \augchow_{\M^F}(y) = (1+y)\augchow_{\M^F\setminus e}(y) +
        \sum_{G \in S_e(\M^F)} y \chow_{\M^F_{G\cup e}}(y) \cdot \augchow_{\M^G}(y).
    \end{equation}
\end{tcolorbox}

\bibliographystyle{amsalpha}
\bibliography{references}

\end{document}
